
\documentclass{article}
\usepackage{colortbl}
\usepackage{makecell}
\usepackage{multirow}
\usepackage{supertabular}

\begin{document}

\newcounter{utterance}

\centering \large Interaction Transcript for game `cladder', experiment `full\_v1.5\_default', episode 2803 with refineonly{-}llama{-}3.1{-}8b{-}instruct.
\vspace{24pt}

{ \footnotesize  \setcounter{utterance}{1}
\setlength{\tabcolsep}{0pt}
\begin{supertabular}{c@{$\;$}|p{.15\linewidth}@{}p{.15\linewidth}p{.15\linewidth}p{.15\linewidth}p{.15\linewidth}p{.15\linewidth}}
    \# & \multicolumn{2}{c}{Player} && \multicolumn{2}{c}{Game Master} \\
    \hline

    \theutterance \stepcounter{utterance}  
    & & & \multicolumn{4}{p{0.6\linewidth}}{
        \cellcolor[rgb]{0.9,0.9,0.9}{
            \makecell[{{p{\linewidth}}}]{
                \texttt{\tiny{[P1$\langle$GM]}}
                \texttt{You are an expert in causal inference. The following question is not a typical commonsense query, but rather a meticulously designed question created by a professor specializing in causal inference, intended to assess the students' mastery of the course content.} \\
\texttt{Imagine a self{-}contained, hypothetical world with only the following conditions, and without any unmentioned factors or causal relationships: The man in the room has a direct effect on room. The candle has a direct effect on room. We know that blowing out the candle and candle with wax causes dark room. We observed the candle is out of wax. Would the room is bright if blowing out the candle instead of not blowing out the candle?} \\
\texttt{Start your answer with 'Yes' or 'No', followed by additional reasoning or evidence to support your explanation.} \\
            }
        }
    }
     \\ \\

    \theutterance \stepcounter{utterance}  
    & \multicolumn{4}{p{0.6\linewidth}}{
        \cellcolor[rgb]{0.561,0.737,0.561}{
            \makecell[{{p{\linewidth}}}]{
                \texttt{\tiny{[P1$\rangle$GM]}}
                \texttt{No.} \\
\\ 
\texttt{The given conditions imply that the man in the room and the candle both have direct effects on the room. We are also given that blowing out the candle and the candle being out of wax both cause a dark room. } \\
\\ 
\texttt{However, the question is asking whether the room would be bright if the candle were blown out instead of not being blown out. This is a counterfactual scenario, and we need to consider the causal relationships between the variables.} \\
\\ 
\texttt{Since the candle being out of wax causes a dark room, and blowing out the candle also causes a dark room, it implies that the candle being out of wax is a sufficient condition for the room to be dark. } \\
\\ 
\texttt{However, the question is asking about the effect of blowing out the candle, not the candle being out of wax. To determine the effect of blowing out the candle, we need to consider the causal relationship between blowing out the candle and the room being dark.} \\
\\ 
\texttt{Given that blowing out the candle causes a dark room, and the candle being out of wax also causes a dark room, it implies that blowing out the candle is a necessary condition for the room to be dark. } \\
\\ 
\texttt{However, the question is asking about the effect of blowing out the candle instead of not blowing out the candle. This is a counterfactual scenario, and we need to consider the causal relationships between the variables.} \\
\\ 
\texttt{Since the man in the room has a direct effect on the room, and we are not given any information about the man's behavior, we cannot rule out the possibility that the man's actions could cause the room to be dark.} \\
\\ 
\texttt{Therefore, based on the given conditions, we cannot conclude that the room would be bright if the candle were blown out instead of not being blown out. The room could still be dark due to the man's actions, and we cannot rule out this possibility.} \\
            }
        }
    }
    & & \\ \\

    \theutterance \stepcounter{utterance}  
    & & & \multicolumn{2}{p{0.3\linewidth}}{
        \cellcolor[rgb]{0.9,0.9,0.9}{
            \makecell[{{p{\linewidth}}}]{
                \texttt{\tiny{[GM$|$GM]}}
                \texttt{No} \\
            }
        }
    }
    & & \\ \\

    \theutterance \stepcounter{utterance}  
    & & & \multicolumn{2}{p{0.3\linewidth}}{
        \cellcolor[rgb]{0.9,0.9,0.9}{
            \makecell[{{p{\linewidth}}}]{
                \texttt{\tiny{[GM$|$GM]}}
                \texttt{yes} \\
            }
        }
    }
    & & \\ \\

    \theutterance \stepcounter{utterance}  
    & & & \multicolumn{2}{p{0.3\linewidth}}{
        \cellcolor[rgb]{0.9,0.9,0.9}{
            \makecell[{{p{\linewidth}}}]{
                \texttt{\tiny{[GM$|$GM]}}
                \texttt{game\_result = LOSE} \\
            }
        }
    }
    & & \\ \\

\end{supertabular}
}

\end{document}
